\section{Groth16 and Pairing-Based SNARKs}
\label{sec:groth16}

The previous chapters developed the Plonk route to SNARKs: computation is represented as a trace, local constraints become polynomial identities, and wiring is enforced by a permutation argument. Groth16 follows a different but equally central route. It starts from a rank-1 constraint system (R1CS), converts that system into a quadratic arithmetic program (QAP), and proves the resulting divisibility relation with a preprocessing pairing argument \cite{GennaroGentryParnoRaykova2013,ParnoHowellGentryRaykova2013,Groth2016}.

For a fixed circuit, Groth16 has perfect completeness, perfect zero knowledge, a proof consisting of two elements of $\G_1$ and one element of $\G_2$, and a verifier whose pairing cost is independent of the circuit size. Its succinctness comes at the price of a circuit-specific structured reference string (SRS): the setup randomness must be sampled correctly and then destroyed.

\subsection{The arithmetization pipeline}

The Groth16 route has the following shape:
\[
\boxed{
\text{arithmetic circuit}
\longrightarrow
\text{R1CS matrices}
\longrightarrow
\text{QAP divisibility}
\longrightarrow
\text{pairing equation}
}.
\]
The conceptual transformation is
\[
\begin{array}{c}
\text{the circuit is evaluated correctly}\\[0.2em]
\Longleftrightarrow\quad \text{all rank-1 constraints are satisfied}\\[0.2em]
\Longleftrightarrow\quad P_z(r)=0\text{ for every constraint point }r\in H\\[0.2em]
\Longleftrightarrow\quad Z_H(X)\mid P_z(X).
\end{array}
\]
Thus many gate equations become one polynomial identity.

\subsection{From algebraic circuits to R1CS}

Let $\F=\F_p$ be a prime field. A satisfying assignment is written
\[
z=(a_0,a_1,\ldots,a_m)\in\F^{m+1},
\qquad a_0=1.
\]
The public coordinates are $(a_0,\ldots,a_\ell)$, and the private witness coordinates are $(a_{\ell+1},\ldots,a_m)$. An R1CS instance with $n$ constraints consists of matrices
\[
A,B,C\in\F^{n\times(m+1)}.
\]
Its $q$th constraint is
\[
\left(\sum_{i=0}^{m}A_{q,i}a_i\right)
\left(\sum_{i=0}^{m}B_{q,i}a_i\right)
=
\sum_{i=0}^{m}C_{q,i}a_i.
\]
Equivalently,
\[
(Az)\circ(Bz)=Cz,
\]
where $\circ$ denotes component-wise multiplication. Addition gates fit this form by multiplying a linear form by the constant wire $a_0=1$; multiplication gates already have the required rank-1 form. Appendix~\ref{app:r1cs-qap-example} gives a color-coded example that tracks every matrix column into the corresponding QAP polynomial.

\subsection{From R1CS to QAP}

Choose distinct constraint points
\[
H=\{r_1,\ldots,r_n\}\subset\F,
\qquad
t(X)=Z_H(X)=\prod_{q=1}^{n}(X-r_q).
\]
For every variable index $i$, interpolate polynomials $u_i,v_i,w_i\in\F[X]$ of degree at most $n-1$ satisfying
\[
u_i(r_q)=A_{q,i},\qquad
v_i(r_q)=B_{q,i},\qquad
w_i(r_q)=C_{q,i}.
\]
The $i$th columns of $A,B,C$ therefore become $u_i,v_i,w_i$. Aggregate them using the assignment:
\[
U(X)=\sum_{i=0}^{m}a_i u_i(X),\qquad
V(X)=\sum_{i=0}^{m}a_i v_i(X),\qquad
W(X)=\sum_{i=0}^{m}a_i w_i(X).
\]
At $r_q$ these polynomials evaluate to the three linear forms in the $q$th R1CS row. Hence all constraints hold if and only if
\[
P_z(r_q)=U(r_q)V(r_q)-W(r_q)=0
\qquad(q=1,\ldots,n).
\]
By the vanishing-polynomial criterion, this is equivalent to the existence of a polynomial $h(X)$ such that
\begin{tcolorbox}
\[
U(X)V(X)-W(X)=h(X)t(X).
\]
\end{tcolorbox}
Since $\deg U,\deg V,\deg W\le n-1$, one may take $\deg h\le n-2$. The QAP relation compresses all $n$ R1CS equations into one divisibility claim.

\subsection{Pairing notation}

Let $\G_1,\G_2,\G_T$ be cyclic groups of prime order $p$, and let
\[
e:\G_1\times\G_2\longrightarrow\G_T
\]
be a non-degenerate asymmetric bilinear pairing. Write
\[
[x]_1=g_1^x,\qquad [x]_2=g_2^x,
\qquad [x]_T=e(g_1,g_2)^x.
\]
Group operations in the source groups add exponents, while the pairing multiplies them:
\[
[x]_j+[y]_j=[x+y]_j,qquad
c[x]_j=[cx]_j,qquad
e([x]_1,[y]_2)=[xy]_T.
\]
The additive notation on $\G_1$ and $\G_2$ is only typographical; the target-group equation below is multiplicative.

\subsection{The structured reference string}

Setup samples
\[
(\alpha,\beta,\gamma,\delta,\tau)\xleftarrow{\$}\F^{5},
\qquad \gamma,\delta\ne 0.
\]
For each public coordinate $0\le i\le\ell$, define
\[
K_i(\tau)=
\frac{\beta u_i(\tau)+\alpha v_i(\tau)+w_i(\tau)}{\gamma},
\]
and for each private coordinate $\ell<i\le m$, define
\[
L_i(\tau)=
\frac{\beta u_i(\tau)+\alpha v_i(\tau)+w_i(\tau)}{\delta}.
\]
The proving key contains the source-group encodings needed to form
$[\alpha+U(\tau)]_1$, $[\beta+V(\tau)]_2$, the private-input linear combination $\sum_{i>\ell}a_i[L_i(\tau)]_1$, and
\[
\left\{\left[\frac{\tau^j t(\tau)}{\delta}\right]_1\right\}_{j=0}^{n-2}.
\]
Concretely this includes suitable encodings of the $u_i(\tau)$ values in $\G_1$, the $v_i(\tau)$ values in $\G_2$ (and the auxiliary encodings required to form the cross terms in $C$), together with $[\delta]_1$ and $[\delta]_2$. The verification key contains
\[
[\alpha]_1,quad [\beta]_2,quad [\gamma]_2,quad [\delta]_2,
\quad \{[K_i(\tau)]_1\}_{i=0}^{\ell}.
\]
The exact published tuple is circuit dependent. Security requires that no party retain the toxic waste $(\alpha,\beta,\gamma,\delta,\tau)$ after setup.

The scalars have separate algebraic roles. The hidden point $\tau$ binds the proof to one QAP identity; $\alpha$ and $\beta$ tie the two multiplicands to the same assignment; $\gamma$ isolates the verifier-controlled public-input contribution; and $\delta$ isolates the witness-dependent part and supports zero-knowledge randomization.

\subsection{Proving}

Given a satisfying assignment and the quotient $h(X)$, the prover samples independent $r,s\xleftarrow{\$}\F$ and forms the exponent-level quantities
\begin{align*}
A
  &=\alpha+U(\tau)+r\delta,\\
B
  &=\beta+V(\tau)+s\delta,\\
C
  &=\frac{\displaystyle
      \sum_{i=\ell+1}^{m}a_i
      \bigl(\beta u_i(\tau)+\alpha v_i(\tau)+w_i(\tau)\bigr)
      +h(\tau)t(\tau)}{\delta}
      +sA+rB-rs\delta.
\end{align*}
The actual proof exposes only group encodings:
\[
\pi=([A]_1,[B]_2,[C]_1)\in\G_1\times\G_2\times\G_1.
\]
Every term is computed as a linear combination of proving-key elements; the prover never learns the setup scalars.

\subsection{Verification and completeness}

From the public coordinates, the verifier computes
\[
[K_{\mathrm{pub}}]_1
=
\sum_{i=0}^{\ell}a_i[K_i(\tau)]_1.
\]
It accepts exactly when
\begin{tcolorbox}
\[
e([A]_1,[B]_2)
\stackrel{?}{=}
e([\alpha]_1,[\beta]_2)\,
e([K_{\mathrm{pub}}]_1,[\gamma]_2)\,
e([C]_1,[\delta]_2).
\]
\end{tcolorbox}
The first pairing may be precomputed as $[\alpha\beta]_T$. Thus online verification uses one pairing-product equation with three pairings, plus a public-input multiscalar multiplication.

To prove completeness, expand the left exponent:
\begin{align*}
AB
&=(\alpha+U+r\delta)(\beta+V+s\delta)\\
&=\alpha\beta+\beta U+\alpha V+UV
  +s\delta(\alpha+U)+r\delta(\beta+V)+rs\delta^2,
\end{align*}
where all polynomials are evaluated at $\tau$. The QAP identity $UV=W+ht$ gives
\[
\beta U+\alpha V+UV
=\sum_{i=0}^{m}a_i(\beta u_i+\alpha v_i+w_i)+ht.
\]
On the other hand,
\begin{align*}
\gamma K_{\mathrm{pub}}+\delta C
&=\sum_{i=0}^{m}a_i(\beta u_i+\alpha v_i+w_i)+ht\\
&\quad+s\delta A+r\delta B-rs\delta^2.
\end{align*}
Substituting $A$ and $B$ into the last line yields precisely all remaining terms of $AB$. Therefore
\[
AB=\alpha\beta+\gamma K_{\mathrm{pub}}+\delta C,
\]
and bilinearity gives the verifier's pairing equation.

\subsection{Zero knowledge and knowledge soundness}

For fixed $\delta\ne0$, the randomizers make $A$ and $B$ uniform over $\F$: translating a uniform field element by $\alpha+U(\tau)$ or $\beta+V(\tau)$ preserves uniformity. Once $A$ and $B$ are fixed, the verification equation determines
\[
C=\frac{AB-\alpha\beta-\gamma K_{\mathrm{pub}}}{\delta}.
\]
A simulator holding the setup trapdoor can therefore sample uniform $A,B$ and compute this $C$ without the witness. The resulting distribution is identical to that of an honest proof, which gives perfect zero knowledge for the original construction.

Soundness is subtler. Informally, the SRS permits the prover to form only prescribed linear combinations in the exponents. The $\alpha$ and $\beta$ terms force the three proof elements to use a consistent assignment; the $\gamma/\delta$ separation prevents private coefficients from being substituted for verifier-controlled public inputs; and the hidden evaluation point makes a false polynomial identity hold only with negligible probability. Groth's analysis formalizes extraction against generic bilinear-group adversaries and establishes an argument of knowledge in that model \cite{Groth2016}. This assumption boundary should not be confused with an information-theoretic proof of knowledge in the standard model.

\subsection{Protocol summary}

\begin{tcolorbox}[enhanced,colback=blue!3!white,colframe=blue!45!black,
                  arc=2.5pt,boxrule=1pt,left=6pt,right=6pt,top=5pt,bottom=5pt]
\small
\textbf{Setup.} Compile the fixed circuit to a QAP, sample
$(\alpha,\beta,\gamma,\delta,\tau)$, publish the proving and verification keys, and destroy the trapdoor.

\medskip
\textbf{Prove.} For a satisfying assignment, compute $h$, sample $r,s$, and output
$\pi=([A]_1,[B]_2,[C]_1)$.

\medskip
\textbf{Verify.} Form $[K_{\mathrm{pub}}]_1$ from the public input and check the single pairing-product equation above.

\medskip
\textbf{Cost.} The proof contains three group elements. Verification uses three online pairings when $e([\alpha]_1,[\beta]_2)$ is precomputed; prover and key sizes scale with the circuit.
\end{tcolorbox}
